\documentclass[11pt]{amsart}

\usepackage[T1]{fontenc}
\usepackage{lmodern}
\usepackage{microtype}
\usepackage{amsmath,amssymb,amsthm}
\usepackage[hidelinks]{hyperref}


\hypersetup{
  pdftitle={An Exhaustive Census for the n=20 Case of the Aouchiche--Caporossi--Hansen Tricyclic-Energy Conjecture}
}

\newtheorem{conjecture}{Conjecture}
\newtheorem{result}{Computational Result}

\newcommand{\energy}{\mathcal E}
\newcommand{\cenergy}{\widehat{\mathcal E}}
\newcommand{\target}{P_{20}^{6,6,6}}

\title[The $n=20$ tricyclic-energy census]
{An Exhaustive Census for the \texorpdfstring{$n=20$}{n=20} Case of the\\
Aouchiche--Caporossi--Hansen Tricyclic-Energy Conjecture}

\date{}

\subjclass[2020]{05C50, 05C35, 68R10}
\keywords{graph energy, tricyclic graph, exhaustive enumeration, nauty}

\begin{document}

\begin{abstract}
Aouchiche, Caporossi, and Hansen conjectured that every tricyclic graph
$G$ on $n=20$ or $n\ge 22$ vertices satisfies
$\energy(G)\le \energy(P_n^{6,6,6})$, with equality only for
$G\cong P_n^{6,6,6}$. We report an exhaustive binary64 census for the
finite case $n=20$. The census contains $3\,288\,208\,176$
isomorphism classes of connected simple graphs with $20$ vertices and
$22$ edges. Its unique computed maximizer is $\target$; the computed
runner-up is separated from it by approximately $1.3061\times10^{-2}$.
Exact characteristic polynomials for both graphs are given. Because
the global comparison was performed in binary64 arithmetic without a
uniform interval certificate, the result is stated as an exhaustive
numerical verification, not as a proof of the conjecture.
\end{abstract}

\maketitle

\section{The finite case}

For a graph $G$ with adjacency eigenvalues
$\lambda_1,\ldots,\lambda_n$, its energy is
\[
  \energy(G)=\sum_{i=1}^n |\lambda_i|.
\]
A connected simple graph is tricyclic when
$|E(G)|=|V(G)|+2$. For $n\ge20$, let $P_n^{6,6,6}$ be the graph
obtained from three copies of $C_6$ and a path $P_{n-18}$ by adding a
single edge between each of two copies of $C_6$ and one endpoint of the
path, and a single edge between the third copy of $C_6$ and the other
endpoint of the path. It has $n$ vertices and $n+2$ edges; no vertices
are identified.

\begin{conjecture}[Aouchiche--Caporossi--Hansen \cite{ACH}, in the form
restated in \cite{LSWL}]
If $G$ is tricyclic on $n=20$ or $n\ge22$ vertices, then
\[
  \energy(G)\le \energy(P_n^{6,6,6}),
\]
with equality if and only if $G\cong P_n^{6,6,6}$.
\end{conjecture}

Let $\cenergy(G)$ denote the energy returned by the binary64
eigensolver used in the census.

\begin{result}\label{res:main}
Among all $3\,288\,208\,176$ isomorphism classes of connected simple
graphs on $20$ vertices and $22$ edges, the reported census found
$\target$ to be the unique maximizer of $\cenergy$. The two largest
computed values are approximately
\[
  27.2903786151833
  \qquad\text{and}\qquad
  27.2773174852287.
\]
\end{result}

For an explicit model of $\target$, take $6$-cycles, in their natural
cyclic orders, on $\{0,\ldots,5\}$, $\{6,\ldots,11\}$, and
$\{12,\ldots,17\}$, and add
\[
  \{0,18\},\qquad
  \{6,18\},\qquad
  \{18,19\},\qquad
  \{12,19\}.
\]
Its characteristic polynomial is
\begin{equation}\label{eq:target-poly}
  \phi_{\target}(x)
  =(x^2-4)(x^2-1)^4
   (x^4-7x^2+8)
   (x^6-7x^4+12x^2-2),
\end{equation}
and hence
\[
  \energy(\target)
  =27.2903786151832973536\ldots .
\]

The runner-up record is isomorphic to the graph
\[
  Q=\target-\{6,18\}+\{5,10\}.
\]
It is not isomorphic to $\target$; for example, $Q$ has six
articulation vertices, whereas $\target$ has five. Its characteristic
polynomial is
\begin{align}\label{eq:runner-poly}
  \phi_Q(x)
  &=(x^2-1)^2\bigl(
  x^{16}-20x^{14}+162x^{12}-686x^{10}+1638x^8 \notag\\
  &\hspace{29mm}{}-2219x^6+1621x^4-557x^2+64\bigr),
\end{align}
so that
\[
  \energy(Q)
  =27.2773174852287231601\ldots
\]
and
\[
  \energy(\target)-\energy(Q)
  =0.0130611299545741935\ldots .
\]

\section{Census}

The program \texttt{geng} from \emph{nauty/Traces} \cite{MP} generates
one representative of each graph isomorphism class. The census was
split into $997$ disjoint shards using
\[
  \texttt{geng -q -c 20 22:22 r/997},
  \qquad 0\le r<997.
\]
The shard counts sum to
\[
  3\,288\,208\,176,
\]
which agrees with the independently tabulated value in OEIS
A001436 \cite{OEIS}. A second count-only run with $512$ shards gave
the same total.

For each graph6 record, the adjacency eigenvalues were computed by
Householder tridiagonalization followed by implicit QL iteration, and
their absolute values were summed. Besides the two threshold counters
displayed next, the scan retained the two largest computed values and
one record attaining each; these are the source of the second value in
Result \ref{res:main} and of the runner-up $Q$. The scan returned
\[
\begin{aligned}
  \#\{G:\cenergy(G)>\cenergy(\target)+10^{-7}\}&=0,\\
  \#\{G:|\cenergy(G)-\cenergy(\target)|\le10^{-7}\}&=1.
\end{aligned}
\]
The unique record in the second set was verified to be isomorphic to
the explicit graph above, and its exact characteristic polynomial
agrees with \eqref{eq:target-poly}. On a comparison sample of
$20\,000$ generated graphs, a second symmetric eigensolver differed
in energy by at most $5.15\times10^{-13}$.

\section{Status and scope}

The observed gap is approximately $1.31\times10^{-2}$, far larger
than the discrepancies observed in the solver comparison. This is
strong numerical evidence for the $n=20$ case. It is not, however, a
uniform error bound over all $3.288$ billion comparisons. Result
\ref{res:main} is therefore a statement about the exhaustive
floating-point census, not an exact-arithmetic theorem.

A proof-level computer-assisted resolution requires a reproducible
archive containing the source, software and build information,
per-shard record accounting, and a validated error bound or interval
certificate sufficient to separate every competitor from $\target$.
What is missing is not margin. Every graph in the census satisfies
$\|A\|_2^2\le\sum_i\lambda_i^2=2|E|=44$, hence $\|A\|_2<7$, and the
standard backward-stability analysis of Householder tridiagonalization
followed by QL iteration bounds the error in each computed eigenvalue
by $c(n)\,\varepsilon\,\|A\|_2$, with $\varepsilon=2^{-53}$ and $c(n)$
a modest polynomial in $n=20$; summed over $20$ eigenvalues, any such
bound stays many orders of magnitude below the observed gap. The
shortfall to a proof therefore lies not in the size of that gap but in
two places no margin can cover: the completeness of the generation
performed by \texttt{geng}, and the applicability of that error model
to the code actually executed on each of the $3.288$ billion inputs.

The present computation concerns only $n=20$ and makes no claim for
$n\ge22$. The maximizer itself is not new here: $\target$ is the graph
already named in the conjecture, so the census confirms a candidate
that was on record beforehand rather than discovering one. The
conjecture in turn came out of heuristic search, in its case the
AutoGraphiX experiments of \cite{ACH}; earlier searches of the same
non-exhaustive kind had been run on related classes, among them the
Monte Carlo-type computer experiment of Gutman and Vidovi\'c on
molecular $(n,m)$-graphs, which recorded the largest-energy graphs the
search found -- not provably maximal ones -- and offered a tentative
guess at the form of the maximal-energy molecular graphs for even $n$
\cite{GV}. What the present census adds is exhaustiveness over all
$3\,288\,208\,176$ isomorphism classes, not the identification of the
maximizer. Previous work proves
the conjecture inside $G(n;a,b,k)$, the connected bipartite tricyclic
graphs on $n\ge20$ vertices with three vertex-disjoint cycles $C_a$,
$C_b$, $C_k$, for most but not all of that class: nine families of
exceptions remain in \cite{LSWL}, and \cite{LMW} continues the analysis
in the same class. Accordingly, this note records an exhaustive
numerical verification of the $n=20$ case; it does not settle the
conjecture.

\section{Exact verification}

Everything this paper asserts about the two exhibited graphs, together
with the census total itself, can be recomputed from the data printed
above in exact integer or rational arithmetic, with no floating-point
eigensolver and no graph-isomorphism package. That recomputation has
been carried out by a second program, written from the definitions and
run independently of the one that produced the census.
It uses its language's standard library only -- no \texttt{nauty}, no
numerical library -- and every decision it takes is made exactly. It
reports $64$ checks and all $64$ pass, but they are not all tests of
this paper. A sizeable part of the count tests the program itself:
that its subclass enumerator is complete and sound, verified
exhaustively against brute force at $n\le10$, and that the auxiliary
count tables it builds along the way reproduce independently tabulated
integer sequences. Those checks license the rest rather than
corroborating anything printed here. What the rest do corroborate, and
what a referee can accordingly redo from this paper alone, is the
following: the explicit model of $\target$ and its
tricyclicity, bipartiteness, degree sequence and three vertex-disjoint
$6$-cycles; the articulation counts $5$ and $6$, and with them the
non-isomorphism of $\target$ and $Q$, established twice over, by
invariants and by an exhaustive backtracking isomorphism search; the
characteristic polynomials
\eqref{eq:target-poly} and \eqref{eq:runner-poly}, each recomputed from
the adjacency matrix and then confirmed a second time by an independent
determinant routine; the energies of $\target$ and $Q$ and their
difference, as rigorous rational enclosures matching every digit
printed here; and the census total $3\,288\,208\,176$, re-derived by
Burnside--P\'olya counting with the multiset logarithm, without
\texttt{geng}. The program and its transcript are retained in an
auxiliary archive, available on request; no claim made here depends on
consulting it, since every check just listed can be reproduced from the
data printed in this paper.

What this recomputation does not do should be stated plainly, since
agreement does not extend to it. It does not repeat the energy scan over all
$3\,288\,208\,176$ classes, which is thousands of CPU-hours. In place
of that it scans exhaustively, with rigorous enclosures, the $8157$
isomorphism classes of connected $20$-vertex $22$-edge graphs all of
whose degrees lie in $\{2,3\}$ -- about $2.5\times10^{-6}$ of the
census, containing both $\target$ and $Q$, and over which $\target$ is
again the unique maximizer and $Q$ the runner-up -- together with the
complete one-edge-swap neighbourhood of $\target$ ($3363$ connected
graphs, most of them outside that class), in which no graph other than
a copy of $\target$ has energy at least $\energy(\target)$. The
remaining classes of the census are not rescanned. Nor does it
reproduce the shard bookkeeping of the $997$- and $512$-shard runs, or
the two binary64 figures of Section 2 -- the counts $0$ and $1$ at the
$10^{-7}$ threshold and the $5.15\times10^{-13}$ solver agreement --
since it replaces floating-point arithmetic by exact enclosures instead
of re-measuring it. The conjecture for $n\ge22$ is outside the scope of
both computations. The exact verification agrees with everything it
checks; it neither extends the census nor turns it into a proof.

\begin{thebibliography}{9}

\bibitem{ACH}
M.~Aouchiche, G.~Caporossi, and P.~Hansen,
\emph{Open problems on graph eigenvalues studied with AutoGraphiX},
EURO J. Comput. Optim. \textbf{1} (2013), 181--199.
\href{https://doi.org/10.1007/s13675-012-0001-9}
{doi:10.1007/s13675-012-0001-9}.

\bibitem{GV}
I.~Gutman and D.~Vidovi\'c,
\emph{Quest for molecular graphs with maximal energy: a computer
experiment},
J. Chem. Inf. Comput. Sci. \textbf{41} (2001), 1002--1005.
\href{https://doi.org/10.1021/ci000164z}
{doi:10.1021/ci000164z}.

\bibitem{LSWL}
X.~Li, Y.~Shi, M.~Wei, and J.~Li,
\emph{On a conjecture about tricyclic graphs with maximal energy},
MATCH Commun. Math. Comput. Chem. \textbf{72} (2014), 183--214.
\href{https://arxiv.org/abs/1312.0204}{arXiv:1312.0204}.

\bibitem{LMW}
X.~Li, Y.~Mao, and M.~Wei,
\emph{More on a conjecture about tricyclic graphs with maximal energy},
MATCH Commun. Math. Comput. Chem. \textbf{73} (2015), 11--26.

\bibitem{MP}
B.~D. McKay and A.~Piperno,
\emph{Practical graph isomorphism, II},
J. Symbolic Comput. \textbf{60} (2014), 94--112.
\href{https://doi.org/10.1016/j.jsc.2013.09.003}
{doi:10.1016/j.jsc.2013.09.003}.

\bibitem{OEIS}
OEIS Foundation Inc.,
\emph{The On-Line Encyclopedia of Integer Sequences},
Sequence A001436.
\url{https://oeis.org/A001436}.

\end{thebibliography}

\end{document}